\documentclass[11pt,a4paper]{article}
\usepackage[utf8]{inputenc}
\usepackage[T1]{fontenc}
\usepackage{amsmath,amssymb,amsthm}
\usepackage{mathtools}
\usepackage{physics}
\usepackage{hyperref}
\usepackage{cleveref}
\usepackage{booktabs}
\usepackage{enumitem}
\usepackage{geometry}
\usepackage{xcolor}
\usepackage{graphicx}
\usepackage{longtable}
\usepackage{caption}

\geometry{margin=2.5cm}

\hypersetup{
  colorlinks=true,
  linkcolor=blue!60!black,
  citecolor=green!50!black,
  urlcolor=blue!70!black,
  pdfauthor={David Alfyorov},
  pdftitle={MR-9: Black Hole Singularity Resolution in the Spectral Action},
}

\newtheorem{theorem}{Theorem}[section]
\newtheorem{proposition}[theorem]{Proposition}
\newtheorem{lemma}[theorem]{Lemma}
\newtheorem{corollary}[theorem]{Corollary}
\theoremstyle{definition}
\newtheorem{definition}[theorem]{Definition}
\theoremstyle{remark}
\newtheorem{remark}[theorem]{Remark}

\definecolor{proven}{rgb}{0,0.5,0}
\definecolor{near}{rgb}{0.8,0.4,0}
\definecolor{distant}{rgb}{0.8,0,0}

\newcommand{\Lam}{\Lambda}
\newcommand{\PiTT}{\Pi_{\mathrm{TT}}}
\newcommand{\Pis}{\Pi_{\mathrm{s}}}
\newcommand{\rs}{r_{\mathrm{s}}}
\newcommand{\rNL}{r_{\mathrm{NL}}}
\newcommand{\Msun}{M_{\odot}}

\title{MR-9: Black Hole Singularity Resolution\\in the Spectral Causal Theory}
\author{David Alfyorov}
\date{April 2026}

\begin{document}
\maketitle

\begin{abstract}
We investigate black hole singularity resolution in the spectral causal theory (SCT),
a candidate framework for UV-complete quantum gravity derived from the spectral action
principle. The linearized SCT field equations produce a modified Schwarzschild metric
with potential $V(r)/V_N(r) = 1 - \tfrac{4}{3}e^{-m_2 r} + \tfrac{1}{3}e^{-m_0 r}$,
where $V(0) = 0$ and the Kretschner scalar is softened from $K \sim 1/r^6$ to $K \sim 1/r^4$.
The singularity is \emph{softened but not resolved} on the linearized level.
We analyze geodesic completeness, Raychaudhuri defocusing, and the nonlinearity radius
$\rNL \sim 1/\Lam$, below which the linearized approximation breaks down.
The Koshelev--Tokareva theorem (2024) implies that for entire functions of order $\geq 1/2$
(SCT has order~1), the singular Schwarzschild metric requires infinite total mass and
is therefore forbidden at the nonperturbative level.
We compare SCT with infinite-derivative gravity (IDG), Stelle gravity, and the
Modesto super-renormalizable program, identifying key discriminating properties.
The tangential null energy condition is violated for $r < \rNL$, providing the
necessary mechanism for Penrose theorem evasion.
Full singularity resolution remains conditional on the explicit evaluation of the
Weyl-sector correction $\Theta^{(C)}_{\mu\nu}$ on curved backgrounds (Gap~G1, OP-01).
\end{abstract}

\tableofcontents

% ===================================================================
\section{Introduction}
% ===================================================================

The black hole singularity problem is one of the central challenges in quantum gravity.
General relativity predicts that gravitational collapse inevitably produces a curvature
singularity at $r = 0$ in the Schwarzschild spacetime, where the Kretschner scalar
$K = R_{\mu\nu\rho\sigma}R^{\mu\nu\rho\sigma} = 48G^2M^2/r^6$ diverges and geodesics
are incomplete~\cite{Penrose1965,HawkingPenrose1970}.

The spectral causal theory (SCT) modifies the Einstein--Hilbert action through nonlocal
form factors derived from the spectral action principle:
\begin{equation}\label{eq:action}
  S = \frac{1}{16\pi G}\int d^4x\,\sqrt{-g}\,
  \bigl[R + \alpha_C\, C_{\mu\nu\rho\sigma}\, F_1(\Box/\Lam^2)\, C^{\mu\nu\rho\sigma}
  + \alpha_R\, R\, F_2(\Box/\Lam^2)\, R\bigr],
\end{equation}
where $\alpha_C = 13/120$ is parameter-free (determined by the Standard Model particle
content), $\Lam$ is the spectral cutoff, and $F_1$, $F_2$ are entire functions of order~1
derived from the master function $\varphi(z) = e^{-z/4}\sqrt{\pi/z}\,\mathrm{erfi}(\sqrt{z}/2)$.

The key question is: does the nonlocal structure of SCT resolve the Schwarzschild singularity?

% ===================================================================
\section{Linearized Analysis}
% ===================================================================

\subsection{Modified Newtonian Potential}

The linearized field equations (NT-4a) yield the modified Newtonian potential in the
Yukawa (pole) approximation:
\begin{equation}\label{eq:potential}
  \frac{V(r)}{V_N(r)} = h(r) = 1 - \frac{4}{3}e^{-m_2 r} + \frac{1}{3}e^{-m_0 r},
\end{equation}
where the effective masses are
\begin{equation}
  m_2 = \Lam\sqrt{\frac{60}{13}} \approx 2.148\,\Lam, \qquad
  m_0 = \Lam\sqrt{6} \approx 2.449\,\Lam \quad (\xi = 0).
\end{equation}

At $r = 0$: $h(0) = 1 - 4/3 + 1/3 = 0$ exactly. The potential vanishes at the origin:
$V(0) = 0$.

\subsection{Modified Metric and L'H\^opital Limit}

The modified Schwarzschild metric function is
\begin{equation}\label{eq:metric}
  f(r) = 1 - \frac{\rs}{r}\,h(r),
\end{equation}
where $\rs = 2GM$ is the Schwarzschild radius. Since $h(0) = 0$, the ratio
$h(r)/r$ requires L'H\^opital's rule:
\begin{equation}\label{eq:lhopital}
  \lim_{r\to 0}\frac{h(r)}{r} = h'(0) = \frac{4m_2}{3} - \frac{m_0}{3}
  \approx 2.048\,\Lam.
\end{equation}

Therefore:
\begin{equation}\label{eq:f_at_zero}
  f(0) = 1 - \rs\,h'(0).
\end{equation}

For astrophysical black holes ($\rs\,\Lam \gg 1$): $f(0) \ll 0$ (deep inside the horizon).
A regular center ($f(0) > 0$) requires $M < M_{\mathrm{crit}} = c^2/(2G\Lam h'(0))$,
which for $\Lam = 2.38\,\mathrm{meV}$ gives $M_{\mathrm{crit}} \approx 1.4 \times 10^{-8}\,M_\odot$
(comparable to the Earth's mass).

\subsection{Kretschner Scalar}

For the metric \eqref{eq:metric}, the Kretschner scalar is
\begin{equation}\label{eq:kretschner}
  K = (f'')^2 + \frac{4(f')^2}{r^2} + \frac{4(f-1)^2}{r^4}.
\end{equation}

With $m(r) = M\,h(r)$ and $h(r) \sim h'(0)\,r$ near $r = 0$:
\begin{equation}
  f(r) \approx 1 - \rs\,h'(0) + O(r), \qquad
  f'(r) \approx -\rs\,h''(0) + O(r),
\end{equation}
the leading Kretschner behavior is
\begin{equation}
  K(r \to 0) \sim \frac{4\bigl(2GM\,h'(0)\bigr)^2}{r^4}.
\end{equation}

This is \emph{softened} from the Schwarzschild $K \sim 48G^2M^2/r^6$ by two powers
of $r$, but \textbf{still diverges}. The singularity is softened but not resolved on
the linearized level.

\subsection{Mass Function Scaling}

The effective enclosed mass $m(r) = M\,h(r)$ behaves as
\begin{equation}
  m(r) \sim M\,h'(0)\,r \quad (r \to 0).
\end{equation}

This is \emph{linear} in $r$, unlike the cubic scaling $m(r) \sim r^3$ characteristic
of a de Sitter core. The absence of a de Sitter core is the root cause of the
unresolved singularity at the linearized level.

% ===================================================================
\section{Nonlinearity Radius and Breakdown of Linearization}
% ===================================================================

\subsection{Definition}

The linearized analysis is valid only for weak gravitational fields. The nonlinearity
radius $\rNL$ marks the scale where the curvature corrections from the form factors
become $O(1)$:
\begin{equation}
  \rNL = \frac{1}{m_2} = \sqrt{\frac{13}{60}}\,\frac{1}{\Lam} \approx \frac{0.465}{\Lam}.
\end{equation}

For $r < \rNL$, the linearized field equations (NT-4a) break down and the full nonlinear
equations (NT-4b) must be used.

\subsection{Kretschner at $r = \rNL$}

At $r = \rNL$, the GR Kretschner scalar for a black hole of mass $M$ is
\begin{equation}
  K_{\mathrm{GR}}(\rNL) = \frac{48G^2M^2}{\rNL^6}
  = 48G^2M^2\Bigl(\frac{60}{13}\Bigr)^3\Lam^6.
\end{equation}

For any astrophysical black hole, $K_{\mathrm{GR}}(\rNL) \gg \Lam^4$, confirming
that the linearized approximation breaks down well before reaching $r = 0$.

\subsection{Koshelev--Tokareva theorem: scope and applicability}

Koshelev and Tokareva~\cite{KoshelevTokareva2024} proved that in quadratic-in-curvature
nonlocal gravity with analytic form factors $F_i(\Box)$, the singular Schwarzschild
metric \emph{as a limit of regular configurations} requires infinite total mass.
Their analysis proceeds via the trace of the full nonlinear equations, integrated
over space, and classifies form factors by their maximal growth in the complex plane
into three cases.

\paragraph{Scope of the no-go.}
The K-T result forbids specifically the \textbf{Schwarzschild $1/r$ branch}---it does
\emph{not} claim that all curvature invariants of the exact solution must be finite.
Our linearized SCT metric has $V(0) = 0$ (no $1/r$ divergence in the potential), with
$m(r) \sim h'(0)\,r \to 0$; the Kretschner divergence $K \sim 1/r^4$ arises from
derivatives, not from the Schwarzschild $m/r^3$ structure.  Therefore our linearized
result is \emph{compatible} with K-T (the Schwarzschild branch is already absent),
rather than contradictory.

\paragraph{Case classification.}
Both standard IDG ($e^{-\Box/M^2}$) and SCT ($\varphi$ of order~1) fall into
\textbf{Case~2}: finite order, growth slower than $\exp(q\,z^{3/2})$.  The strong
no-go (Case~1, requiring infinite mass) applies only to faster-growing form factors
of order ${}>3/2$ or infinite order.  In Case~2, the convergence of the mass integral
is \emph{regularization-dependent}: it depends on the family of smooth metrics used
to approach the singular limit, not on the QFT pole prescription.  Thus K-T does
\emph{not} resolve the question for SCT in either direction.

\paragraph{Fakeon caveat.}
K-T analyze the variational equations of the full action.  The fakeon prescription
modifies the physical classical equations (replacing the Feynman propagator by the
average of retarded and advanced Green functions for the ghost mode), so the
fakeon-projected classical dynamics differs from the naively varied action.
Consequently, K-T cannot be applied directly to the physical SCT dynamics without
repeating their mass-integral analysis for the fakeon-projected equations.

\paragraph{Global order vs.\ Euclidean asymptotics.}
The master function $\varphi(z)$ has order~1 \emph{globally} (growth $\sim e^{|z|/4}$
on the negative real axis), but $\PiTT(z) \to \mathrm{const}$ on the positive real
axis because $\varphi(z) \sim 2/z$ there.  These are compatible: the order is a
\emph{global complex-plane} property controlling K-T convergence, while the propagator
UV behavior on the Euclidean axis determines the linearized source smearing.

\paragraph{Nonlinearity radius.}
For black holes with $\rs \gg 1/\Lam$, the linearized metric $f(r) = 1 - (\rs/r)\,h(r)$
has $|f - 1| = O(1)$ already at $r \sim \rs$, not at $r \sim 1/\Lam$.  The correct
breakdown scale is $r_{\rm NL} \sim \max(\rs,\, 1/\Lam)$; for astrophysical black holes,
the entire interior $r < \rs$ is outside the domain of validity of the linearized theory.

\paragraph{Implication.}
K-T provides \emph{weak} support for singularity resolution in SCT: it excludes the
Schwarzschild $1/r$ branch but does not exclude a milder, non-Schwarzschild singularity
(such as $K \sim 1/r^4$).  The possibility of a softer-than-Schwarzschild but still
singular exact solution remains open.  Full resolution requires either (a) solving the
fakeon-projected nonlinear equations (blocked by Gap~G1), or (b) repeating K-T's
mass-integral analysis for the fakeon-classical equations.

% ===================================================================
\section{Geodesic Completeness}
% ===================================================================

\subsection{Monotonicity of $f(r)$ and horizon structure}

Define $g(r) = h(r) - r\,h'(r)$. Then $g(0) = 0$ and $g'(r) = -r\,h''(r) > 0$
(since $h''(r) < 0$ for all $r > 0$), so $g(r) > 0$ for $r > 0$. Since
$f'(r) = \rs\, g(r)/r^2 > 0$, the metric function $f(r)$ is \textbf{monotonically
increasing} from $f(0) = 1 - \rs\,h'(0)$ to $f(\infty) = 1$. Consequently,
for $\rs\,h'(0) > 1$ (all astrophysical black holes), there is \textbf{exactly one
horizon} where $f = 0$. The center $r = 0$ lies in a $T$-region where $r$ is
timelike, the same causal structure as the Schwarzschild interior.

Note: $M_{\mathrm{crit}}$ defined by $f(0) = 0$ ($\rs\,h'(0) = 1$) is the
\textbf{horizon existence threshold}, not a regularity threshold. A regular
center requires $f(0) = 1$ and $f'(0) = 0$, equivalently $m(r) = O(r^3)$.
Since $m(r) = (\rs/2)\,h'(0)\,r + O(r^2)$ for any $M > 0$, the linearized
SCT metric is \textbf{never regular at $r = 0$}, regardless of the mass.

\subsection{Geodesic incompleteness}

For a general causal geodesic with energy $E$ and angular momentum $L$:
\begin{equation}
  \dot{r}^2 = E^2 - f(r)\Bigl(\kappa + \frac{L^2}{r^2}\Bigr),
  \qquad \kappa = \begin{cases} 1 & \text{timelike,} \\ 0 & \text{null.}\end{cases}
\end{equation}
Near $r = 0$ with $f_0 := f(0) < 0$:
\begin{itemize}
  \item Radial timelike ($L = 0$):
    $\dot{r}^2 \to E^2 - f_0 > 0$, so $\tau \sim r/\sqrt{E^2 - f_0} \to 0$ (finite).
  \item Radial null ($L = 0$):
    $\dot{r}^2 = E^2$, so $\lambda \sim r/E \to 0$ (finite).
  \item Nonzero $L$:
    $\dot{r}^2 \sim |f_0|\,L^2/r^2$, so $\Delta\tau \sim O(r^2)$ (even more convergent).
\end{itemize}
All causal geodesics reach $r = 0$ in finite affine/proper parameter.

The key question is whether $r = 0$ is a regular interior point or a non-extendable
boundary. In Cartesian coordinates $x^i = r\,n^i$, the metric becomes
$g_{ij} = \delta_{ij} + (f^{-1} - 1)\,x_i x_j / r^2$. If $f(0) \neq 1$, the
limit at $r \to 0$ depends on the direction $x_i/r$, so $r = 0$ cannot be added as
a regular interior point. Combined with $K \sim (1 - f_0)^2/r^4 \to \infty$, the
spacetime is \textbf{geodesically incomplete}: geodesics reach $r = 0$ in finite
parameter, and $r = 0$ is not extendable.

\subsection{What nonlinear corrections must achieve}

For geodesic completeness, the nonlinear corrections must replace the linearized
near-center asymptotics $m(r) \sim \alpha\,r$ with $m(r) \sim \beta\,r^3$, yielding
$f(r) = 1 + O(r^2)$ (de Sitter core). This is a \textbf{qualitative restructuring},
not a small perturbation: the entire leading-order behavior at $r = 0$ must change.
Even if a local de Sitter core is achieved, global maximal extendability of the
resulting geometry through all horizons requires separate verification.

\paragraph{Kinematic structure of the regular center.}
In the two-function ansatz $ds^2 = -A(r)\,dt^2 + B(r)\,dr^2 + r^2\,d\Omega^2$,
regularity of $G^t{}_t$ at $r = 0$ requires $b_1 = 0$; regularity of $G^r{}_r$
requires $a_1 = b_1$.  Hence a smooth center admits only even powers:
$A = 1 + a_2 r^2 + \cdots$, $B = 1 + b_2 r^2 + \cdots$.  The mass function
$m(r) = \tfrac{r}{2}(1 - 1/B) = \tfrac{b_2}{2}\,r^3 + O(r^5)$
is \emph{automatically} cubic ($n = 3$).  This is kinematic, not dynamic:
$n = 3$ is a \emph{consequence} of smoothness, not an independent condition.

\paragraph{Branch structure.}
The $n = 1$ result is stable on the Schwarzschild/Yukawa-connected perturbative
branch: no finite order in $\varepsilon = GM\Lam$ can cancel the leading linear term.
A regular $n = 3$ solution, if it exists, must be a \textbf{separate nonperturbative
branch}, not a perturbative correction.  Such branch structure is typical of
Weyl-sector higher-derivative gravity (L\"u--Perkins--Pope--Stelle, 2015).
The fakeon prescription does not help: classicization is itself perturbative and
cannot produce a branch change.

\paragraph{Selection by collapse.}
In quadratic gravity, dust collapse produces horizons and singularities, with the
$C^2$ terms even accelerating singularization.  There is no evidence that collapse
dynamics select a regular branch over the Schwarzschild-connected one in SCT.

% ===================================================================
\section{Raychaudhuri Defocusing}
% ===================================================================

\subsection{Null Energy Condition Analysis}

For the metric $f(r) = 1 - 2Gm(r)/r$ with $m(r) = M\,h(r)$, the effective
stress-energy has:
\begin{itemize}
  \item Radial NEC: $\rho + p_r = 0$ (marginal, identically satisfied).
  \item Tangential NEC: $\rho + p_t = (2m' - r\,m'')/(8\pi r^2)$.
\end{itemize}

The correct tangential NEC quantity (verified against the de~Sitter cross-check
$m(r) = \rho_0 r^3/3$, for which $\rho + p_t = 0$ identically) is
\begin{equation}
  g(r) \equiv 2m'(r) - r\,m''(r) = M\bigl[2h'(r) - r\,h''(r)\bigr].
\end{equation}

Although $h''(r) < 0$ for all $r > 0$ (the mass function is everywhere concave),
the quantity $g(r) > 0$ for all $r \geq 0$:
\begin{equation}
  g(0) = 2h'(0) = 2\Bigl(\frac{4m_2}{3} - \frac{m_0}{3}\Bigr) > 0,
\end{equation}
and $g(r) \to 0^+$ as $r \to \infty$. Numerical verification at 100-digit precision
confirms $g(r) > 0$ for all $r \in (0, 100/\Lam)$.

\textbf{The tangential NEC is NEVER violated in the linearized SCT Yukawa metric.}

\subsection{Penrose Theorem Status}

Since the NEC is satisfied everywhere on the linearized SCT metric, all conditions
of the Penrose singularity theorem~\cite{Penrose1965} are met:
\begin{enumerate}
  \item NEC: satisfied (both radially and tangentially).
  \item Trapped surface: present (horizons exist for $M > M_{\mathrm{crit}}$).
  \item Global hyperbolicity: standard for asymptotically flat spacetimes.
  \item Cauchy surface: exists.
\end{enumerate}

The Penrose theorem therefore \emph{applies} at the linearized level and predicts
geodesic incompleteness. This is fully consistent with the Kretschner result
$K \sim 1/r^4$.

\textbf{Comparison with IDG:} In ghost-free IDG with $\exp(-\Box/M^2)$, the
exponential propagator modification produces a Gaussian smearing of the source
($\rho_{\mathrm{eff}} \sim e^{-r^2 M^2/4}$), which \emph{does} violate the NEC
near $r = 0$, enabling singularity resolution at the linearized level. SCT's
order-1 entire function $\varphi(z) \sim 2/z$ produces $\PiTT \to \mathrm{const}$
in the UV, which is insufficient for linearized NEC violation.

The critical distinction is: IDG has $\Pi \to \infty$ (exponential growth);
SCT has $\PiTT \to \mathrm{const}$ (saturation). Both have entire-function
order~1, but the propagator UV behavior differs qualitatively.

% ===================================================================
\section{Comparison with Other Frameworks}
% ===================================================================

\begin{table}[htbp]
\centering
\caption{Black hole singularity resolution across modified gravity frameworks.}
\label{tab:comparison}
\small
\begin{tabular}{@{}lllll@{}}
\toprule
Property & SCT & IDG & Stelle & Hayward \\
\midrule
Action origin & Spectral & Postulated & Quadratic & Metric ansatz \\
Entire-fn order & 1 & 1 & N/A & N/A \\
$\PiTT$ UV & $\to\,$const & $\to\,\infty$ & $\to\,$const & N/A \\
Source smearing & No (lin.) & Gaussian & Partial & By constr. \\
Core type & No dS (lin.) & de Sitter & No dS & de Sitter \\
$K$ scaling & $1/r^4$ & Finite & $1/r^4$ & Finite \\
$m(r\to 0)$ & $\sim r$ & $\sim r^3$ & $\sim r$ & $\sim r^3$ \\
Geodesic & Incomplete & Complete & Incomplete & Complete \\
Ghost? & Fakeon & None & Massive & N/A \\
Free params & 0 & 1 & 2 & 1 \\
Singularity & Softened & Resolved & Not resolved & Resolved \\
\bottomrule
\end{tabular}
\end{table}

SCT occupies a unique position: it is derived from a deeper principle (spectral action)
with zero free parameters beyond $\Lam$, but its linearized predictions match those of
Stelle gravity (softened but unresolved).

\paragraph{The $n$-classification.}
The near-center behavior $m(r) \sim c\,r^n$ classifies singular vs.\ regular cores:
$n < 3$ (SCT: $n = 1$): $\rho + p_t = c\,n(3-n)\,r^{n-3}/(8\pi) > 0$ (focusing,
singular); $n = 3$ (de Sitter): marginal; $n > 3$: defocusing.  SCT sits deep in the
singular class, two powers of $r$ below the de Sitter threshold.

\paragraph{Smoothing power vs.\ entire-function order.}
What determines $n$ is not the complex-plane order of $\varphi(z)$, but the UV
behavior of the propagator denominator $\PiTT$ on the \emph{physical Euclidean axis}:
$\PiTT \to \mathrm{const}$ gives $n = 1$ (zero extra smoothing---same UV as GR);
$\PiTT \sim k^{2N}$ gives $n = 2N + 1$;
$\PiTT \sim e^{k^2/M^2}$ (IDG) gives $n = 3$ (Gaussian smearing, de Sitter core).
Nicolini's order-$\geq 1/2$ threshold applies to his specific class of cutoffs, not
universally.

\paragraph{Bach-flatness and the $C\Box C$ correction (Gap~G1 partial resolution).}
Schwarzschild is an Einstein space, hence Bach-flat ($B_{\mu\nu} = 0$).  In
\emph{local} $C^2$ gravity, this means Schwarzschild is an exact vacuum solution.

In the nonlocal theory, the first surviving contribution comes from the $n = 1$ term
in the Taylor expansion $\hat{F}_1(z) = 1 + \hat{f}_1 z + O(z^2)$, i.e., from the
$C\Box C$ term in the action.  We have computed (and verified at 100-digit precision):
\begin{equation}
  \hat{f}_1 \equiv \hat{F}_1'(0) = -\frac{307}{182} \approx -1.687.
\end{equation}
On Schwarzschild, the covariant d'Alembertian of the Weyl tensor satisfies the
remarkable algebraic identity (in orthonormal frame):
\begin{equation}\label{eq:box_C}
  \Box\, C_{\hat{a}\hat{b}\hat{c}\hat{d}}
  = -\frac{6M}{r^3}\, C_{\hat{a}\hat{b}\hat{c}\hat{d}},
\end{equation}
so $C\,\Box C \sim C^3 \sim M^3/r^9$.  The first nonzero $\Theta^{(C)}_{\mu\nu}$ on
Schwarzschild therefore scales as $\alpha_C\,\hat{f}_1\,\Lam^{-2}\,M^2/r^8$ and
falls off as $1/r^8$--$1/r^9$.  Explicit $a_6$ (six-derivative spectral action)
computations confirm that Schwarzschild is \emph{not} an exact solution at this order.

However, this correction \emph{adds} a cubic term to the mass function
($m(r) = a_1 r + a_3 r^3 + \cdots$) rather than cancelling the linear term $a_1$.
The UV saturation $\PiTT \to \mathrm{const}$ provides no mechanism to set $a_1 = 0$.

% ===================================================================
\section{Nonperturbative Spectral Action Argument}
% ===================================================================

A potential bypass of Gap~G1 is the nonperturbative spectral action argument:
the spectral action $S = \mathrm{Tr}\,f(D^2/\Lam^2)$ assigns \emph{infinite} action
to the singular Schwarzschild geometry (due to $\int C^2\,\sqrt{g}\,d^4x = \infty$)
and \emph{finite} action to regular geometries (Hayward, Bardeen).

If the path integral weights geometries by $e^{-S}$, the singular geometry has zero
weight and is dynamically excluded. This argument is heuristic but suggestive: it
provides a mechanism for singularity resolution that does not require explicit
knowledge of $\Theta^{(C)}_{\mu\nu}$.

\textbf{Caveat:} This argument relies on the perturbative expansion of the spectral
action. The full nonperturbative spectral action on a singular background is not
well-defined (the Dirac operator is not essentially self-adjoint at $r = 0$), which
may itself be interpreted as forbidding the singular geometry.

% ===================================================================
\section{Verdict and Open Questions}
% ===================================================================

\subsection{Summary of Results}

\begin{enumerate}
  \item \textbf{Linearized:} $K \sim 1/r^4$ (softened from $1/r^6$, NOT resolved).
    Mass function $m(r) \sim r$ (linear, not cubic). No de Sitter core.
  \item \textbf{Geodesics:} Incomplete on linearized metric (same as Schwarzschild).
  \item \textbf{NEC:} Tangential NEC \emph{satisfied} everywhere (2$m'$ -- $r m'' > 0$).
    Penrose theorem \emph{applies} at linearized level.
  \item \textbf{Nonlinearity:} Linearization breaks down at $\rNL \sim 0.47/\Lam$.
    All singular behavior ($K \sim 1/r^4$) occurs in the unreliable region.
  \item \textbf{K-T theorem:} SCT is Case~2 (same as IDG). K-T forbids
    the Schwarzschild $1/r$ branch (already absent in SCT linearized),
    but does \emph{not} exclude a milder singularity. Fakeon caveat:
    K-T may not apply to fakeon-projected equations.
  \item \textbf{Nonperturbative:} $S[\mathrm{Schw}] = \infty > S[\mathrm{Hayward}] < \infty$.
    Spectral action energetically prefers regularity.
\end{enumerate}

\subsection{Overall Verdict}

\begin{center}
\fbox{\parbox{0.85\textwidth}{
\textbf{SINGULARITY RESOLUTION: NEGATIVE at linearized level; OPEN at nonlinear.}\\[4pt]
Linearized: $K \sim 1/r^4$, $m(r) \sim r$ ($n = 1$, deep in singular class,
two powers below de Sitter $n = 3$). NEC satisfied everywhere; Penrose theorem
applies; all causal geodesics incomplete.  Root cause: $\PiTT \to \mathrm{const}$
provides zero extra UV smoothing.\\[4pt]
Nonlinear: \textbf{unknown}.  Schwarzschild is Bach-flat ($B_{\mu\nu} = 0$), so
the leading Weyl-sector correction vanishes; but nonlocal commutator terms
$[\Box^n, C]$ may push the solution off the Schwarzschild branch.  Whether this
changes $m(r) \sim r \to r^3$ is \emph{precisely} what Gap~G1 (OP-01) must resolve.
K-T forbids only the Schwarzschild $1/r$ branch (already absent here);
a milder singularity is not excluded.  The $n$-classification and UV smoothing
analysis suggest \textbf{pessimism} but do not constitute a proof of persistence.
}}
\end{center}

\subsection{Open Questions}

\begin{enumerate}
  \item \textbf{OP-01 (Gap G1):} Compute $\Theta^{(C)}_{\mu\nu}$ on Schwarzschild.
    This unblocks the full nonlinear analysis.
  \item \textbf{K-T Case~2:} Is SCT's resolution regularization-independent?
    Does the fakeon prescription affect the outcome?
  \item \textbf{Self-consistent solution:} Solve the full NT-4b equations on a
    static spherically symmetric ansatz. Determine if the solution has a de Sitter
    core, bounce, or other regular structure.
  \item \textbf{Nonperturbative spectral action:} Compute
    $S[\mathrm{Hayward}(L)]$ as a function of $L$ and find the minimum.
    This gives the predicted core radius without solving the field equations.
\end{enumerate}

% ===================================================================
\section*{Acknowledgments}
% ===================================================================

I thank Igor Shnyukov for discussions on the nonlinear field equations
and the spectral action on singular backgrounds.

% ===================================================================
\begin{thebibliography}{30}
\bibitem{Penrose1965}
R.~Penrose, ``Gravitational collapse and space-time singularities,''
\textit{Phys. Rev. Lett.} \textbf{14}, 57 (1965).

\bibitem{HawkingPenrose1970}
S.~W.~Hawking and R.~Penrose, ``The singularities of gravitational collapse
and cosmology,'' \textit{Proc. R. Soc. Lond. A} \textbf{314}, 529 (1970).

\bibitem{Stelle1977}
K.~S.~Stelle, ``Renormalization of higher-derivative quantum gravity,''
\textit{Phys. Rev. D} \textbf{16}, 953 (1977).

\bibitem{BGKM2012}
T.~Biswas, E.~Gerwick, T.~Koivisto, and A.~Mazumdar,
``Towards singularity and ghost free theories of gravity,''
\textit{Phys. Rev. Lett.} \textbf{108}, 031101 (2012),
\href{https://arxiv.org/abs/1110.5249}{arXiv:1110.5249}.

\bibitem{Modesto2012}
L.~Modesto, ``Super-renormalizable quantum gravity,''
\textit{Phys. Rev. D} \textbf{86}, 044005 (2012),
\href{https://arxiv.org/abs/1107.2403}{arXiv:1107.2403}.

\bibitem{Nicolini2012}
P.~Nicolini, ``Nonlocal and generalized uncertainty principle black holes,''
\href{https://arxiv.org/abs/1202.2102}{arXiv:1202.2102} (2012).

\bibitem{KoshelevTokareva2024}
A.~S.~Koshelev and A.~Tokareva,
``Non-perturbative quantum gravity denounces singular black holes,''
\textit{Phys. Rev. D} \textbf{111}, 086026 (2025),
\href{https://arxiv.org/abs/2404.07925}{arXiv:2404.07925}.

\bibitem{ConroyMazumdar2017}
A.~Conroy and A.~Mazumdar,
``Newtonian potential and geodesic completeness in infinite derivative gravity,''
\href{https://arxiv.org/abs/1705.02382}{arXiv:1705.02382} (2017).

\bibitem{EdholmKM2016}
J.~Edholm, A.~S.~Koshelev, and A.~Mazumdar,
``Behavior of the Newtonian potential for ghost-free gravity and singularity free gravity,''
\textit{Phys. Rev. D} \textbf{94}, 104033 (2016),
\href{https://arxiv.org/abs/1604.01989}{arXiv:1604.01989}.

\bibitem{Hayward2006}
S.~A.~Hayward, ``Formation and evaporation of nonsingular black holes,''
\textit{Phys. Rev. Lett.} \textbf{96}, 031103 (2006).

\bibitem{Burzilla2023}
N.~Burzilla et al., ``Regular multi-horizon Lee-Wick black holes,''
\href{https://arxiv.org/abs/2308.12810}{arXiv:2308.12810} (2023).

\bibitem{LMR2015}
Y.-D.~Li, L.~Modesto, and L.~Rachwal,
``Exact solutions and spacetime singularities in nonlocal gravity,''
\textit{JHEP} \textbf{12}, 173 (2015),
\href{https://arxiv.org/abs/1506.08619}{arXiv:1506.08619}.

\bibitem{Alfyorov2026a}
D.~Alfyorov, ``Nonlocal one-loop form factors from the spectral action,''
DOI:\href{https://doi.org/10.5281/zenodo.19098042}{10.5281/zenodo.19098042} (2026).

\bibitem{Alfyorov2026b}
D.~Alfyorov, ``Solar system and laboratory tests of the spectral action,''
DOI:\href{https://doi.org/10.5281/zenodo.19098100}{10.5281/zenodo.19098100} (2026).

\bibitem{Alfyorov2026c}
D.~Alfyorov, ``Nonlinear field equations and FLRW cosmology from the spectral action,''
DOI:\href{https://doi.org/10.5281/zenodo.19098027}{10.5281/zenodo.19098027} (2026).

\end{thebibliography}

\end{document}
